% Generated draft: requires mathematical and source review.
\documentclass[11pt]{article}
\usepackage[margin=25mm]{geometry}
\usepackage{fontspec}
\setmainfont{Latin Modern Roman}
\usepackage{amsmath,amssymb,mathtools,unicode-math}
\setmathfont{Latin Modern Math}
\usepackage{enumitem,needspace,xurl,hyperref}
\setlength{\parindent}{0pt}
\setlength{\parskip}{5pt}
\setlength{\emergencystretch}{3em}
\newcommand{\N}{\mathbb{N}}
\newcommand{\Z}{\mathbb{Z}}
\newcommand{\Q}{\mathbb{Q}}
\newcommand{\R}{\mathbb{R}}
\newcommand{\C}{\mathbb{C}}
\newcommand{\F}{\mathbb{F}}
\newcommand{\A}{\mathbb{A}}
\newcommand{\PP}{\mathbb{P}}
\newcommand{\E}{\mathbb{E}}
\newcommand{\eps}{\varepsilon}
\newcommand{\dd}{\,\mathrm{d}}
\begin{document}
\section*{UnsolvedMath ID: 700461}
\textbf{Source problem code:} catalogue entry
\subsection{Definitions and mathematical statement}
Let \(G\) be an arbitrary group. For \(S\subseteq G\) and \(x\in G\), define the left and right translates by
\[
xS=\{xs:s\in S\},\qquad Sx=\{sx:s\in S\}.
\]
Following the supplied statement attributed to C. Chou, say that \(G\) has property \(P\) if, for every finite subset \(F\subseteq G\), there exist a finite subset \(S\subseteq G\) and a subset \(X\subseteq G\) such that
\[
S\supset F,\qquad
\forall x_1,x_2\in X\;\bigl(x_1\ne x_2\Longrightarrow x_1S\cap x_2S=\varnothing\bigr),\qquad
G=\bigcup_{x\in X}Sx.
\]
Here \(X\) is not required to be finite. The notation \(S\supset F\) is retained from the source; the record does not explicitly specify whether this containment must be proper.

\textbf{Question.} Does every group \(G\) have property \(P\)?

\textbf{Unresolved source qualification.} The disjointness condition concerns left translates \(xS\), whereas the covering condition concerns right translates \(Sx\), using the same set \(X\). These conditions are retained exactly as supplied. The record does not explain whether this difference is intentional; no common-sided tiling formulation is substituted.
\subsection{Short English statement}
Does every group admit, around each finite subset, a finite set whose left translates by a suitable set of group elements are pairwise disjoint and whose right translates by those same elements cover the group? The source leaves the strictness of containment unspecified and uses different translation sides for disjointness and covering.
\subsection{Sources}
\begin{itemize}
\item Local UnsolvedMath registry, permanent ID \(700461\), exact source problem code \texttt{catalogue entry}, title \emph{Kourovka Notebook Problem 8.9}. The supplied registry labels this problem \texttt{open}; that label is supplied source metadata, not an independently verified mathematical assessment.
\item \emph{The Kourovka Notebook (Unsolved Problems in Group Theory)}, original problem 8.9, 8th issue (1982); supplied reference edition: 21st issue (2026), arXiv:1401.0300v48. The supplied collection description names E. I. Khukhro and V. D. Mazurov as editors.
\item The registry identifies R. I. Grigorchuk as proposer, with proposed year 1982. The statement itself begins with the attribution ``(C. Chou)''; both attributions are preserved without inferring their precise relationship.
\end{itemize}
\end{document}
